\documentclass[aspectratio=169,10pt]{beamer}

% Shared Beamer setup for Fall 2026 course slides.
% Clean Cal Poly-inspired theme: no custom class files or uncommon packages.
\mode<presentation>{
  \usetheme{default}
  \usefonttheme{professionalfonts}
}

\definecolor{CalPolyGreen}{HTML}{154734}
\definecolor{CalPolyGold}{HTML}{BD8B13}
\definecolor{CalPolySage}{HTML}{7FA28F}
\definecolor{CalPolyMint}{HTML}{DDF1E7}
\definecolor{CalPolySand}{HTML}{A29F86}
\definecolor{CalPolyBeige}{HTML}{EFEEDE}
\definecolor{DeckInk}{HTML}{1F2933}
\definecolor{DeckMuted}{HTML}{5F6B7A}
\definecolor{DeckSoft}{HTML}{F7F7F5}

\setbeamersize{text margin left=0.55in,text margin right=0.55in}
\setbeamercolor{normal text}{fg=DeckInk,bg=white}
\setbeamercolor{structure}{fg=CalPolyGreen}
\setbeamercolor{title}{fg=white}
\setbeamercolor{frametitle}{fg=CalPolyGreen,bg=white}
\setbeamercolor{block title}{bg=CalPolySage,fg=white}
\setbeamercolor{block body}{bg=CalPolyMint,fg=black}
\setbeamercolor{alerted text}{fg=CalPolyGold}

\setbeamerfont{title}{size=\Huge,series=\bfseries}
\setbeamerfont{frametitle}{size=\Large,series=\bfseries}
\setbeamerfont{block title}{size=\normalsize,series=\bfseries}

\setbeamertemplate{navigation symbols}{}
\setbeamertemplate{itemize item}{\color{CalPolyGold}$\blacktriangleright$}
\setbeamertemplate{itemize subitem}{\color{CalPolyGreen}$\bullet$}
\setbeamertemplate{footline}{
  \hfill{\color{DeckMuted}\scriptsize\insertframenumber/\inserttotalframenumber}
  \hspace{0.18in}\vspace{0.12in}
}
\setbeamertemplate{frametitle}{
  \vspace{0.18in}
  \hspace{0.02in}\insertframetitle\par
  \vspace{0.08in}
  \noindent\makebox[\linewidth][l]{%
    {\color{CalPolyGold}\rule{0.55in}{2pt}}%
    {\color{CalPolyGreen}\rule{\dimexpr\linewidth-0.55in\relax}{2pt}}%
  }
}

\newcommand{\CourseNumber}{Course}
\newcommand{\CourseTitle}{Course Title}
\newcommand{\LectureNumber}{0}
\newcommand{\LectureTitle}{Lecture Title}
\newcommand{\LectureDate}{Date}
\newcommand{\CourseUnit}{Unit}

\newcommand{\coursenumber}[1]{\renewcommand{\CourseNumber}{#1}}
\newcommand{\coursetitle}[1]{\renewcommand{\CourseTitle}{#1}}
\newcommand{\lecturenumber}[1]{\renewcommand{\LectureNumber}{#1}}
\newcommand{\lecturetitle}[1]{\renewcommand{\LectureTitle}{#1}}
\newcommand{\lecturedate}[1]{\renewcommand{\LectureDate}{#1}}
\newcommand{\courseunit}[1]{\renewcommand{\CourseUnit}{#1}}

\author{Kirk Duran}
\institute{California Polytechnic State University, San Luis Obispo}
\date{\LectureDate}

\newcommand{\makedecktitle}{
  \begingroup
  \setbeamercolor{background canvas}{bg=CalPolyGreen}
  \begin{frame}[plain]
    \vfill
    \begin{center}
      {\usebeamerfont{title}\color{white}\LectureTitle\par}
      \vspace{0.18in}
      {\Large\color{white}Lecture \LectureNumber\par}
    \end{center}
    \vfill
    \noindent{\color{white}\rule{\linewidth}{1.2pt}}\par
    \vspace{0.08in}
    \noindent\colorbox{CalPolyGold}{\makebox[0.24\linewidth][c]{\color{white}\Large\CourseNumber}}
    \hspace{0.12in}{\color{white}\Large Kirk Duran}
  \end{frame}
  \endgroup
}

\newenvironment{keyidea}{\begin{block}{Key idea}}{\end{block}}
\newenvironment{activity}{\begin{block}{In class}}{\end{block}}
\newenvironment{cpdefinition}{\begin{block}{Definition}}{\end{block}}
\newenvironment{exampleblockcp}{
  \begingroup
  \setbeamercolor{block title}{bg=CalPolySand,fg=white}
  \setbeamercolor{block body}{bg=CalPolyBeige,fg=black}
  \begin{block}{Example}
}{
  \end{block}
  \endgroup
}

\newcommand{\imageplaceholder}[2][0.3in]{%
  \noindent\makebox[\linewidth][c]{%
    \fcolorbox{CalPolySage}{CalPolyBeige}{%
      \parbox[c][#1][c]{0.86\linewidth}{%
        \centering
        {\color{DeckMuted}\scriptsize Image placeholder: }%
        {\color{DeckInk}\scriptsize\ttfamily\detokenize{#2}}%
      }%
    }%
  }%
}

\newcommand{\sectionframe}[1]{
  \begingroup
  \setbeamercolor{background canvas}{bg=CalPolyGreen}
  \begin{frame}[plain]
    \vfill
    {\color{white}\LARGE\bfseries #1\par}
    \vspace{0.18in}
    {\color{CalPolyGold}\rule{0.22\paperwidth}{3pt}}
    {\color{white}\rule{0.62\paperwidth}{3pt}}
    \vfill
  \end{frame}
  \endgroup
}

\newcommand{\placeholdernote}{
  \begin{block}{Placeholder}
    Replace these bullets with the final lecture material, examples, and in-class activities.
  \end{block}
}

\usepackage[T1]{fontenc}
\usepackage{lmodern}
\pdfstringdefDisableCommands{\def\translate#1{#1}}

\graphicspath{{../assets/lecture-30/}}

% Generated classroom visuals live in ../assets/lecture-30/.
\newcommand{\lectureimage}[2][1.75in]{%
  \includegraphics[width=\linewidth,height=#1,keepaspectratio]{#2}%
}

\setbeamersize{text margin left=0.42in,text margin right=0.42in}

\coursenumber{CS 1001}
\coursetitle{Introduction to Computing}
\lecturenumber{30}
\lecturetitle{Searching Algorithms}
\lecturedate{Fall 2026}
\courseunit{12. Algorithms}

\begin{document}

\makedecktitle

\begin{frame}{Today}
  \begin{itemize}
    \item Define an \textbf{algorithm} as a precise procedure for solving a computational problem.
    \item State the search problem independently of any particular Python implementation.
    \item Implement and trace \textbf{linear search} and reason about how much work it performs.
    \item Search lists of objects by examining selected attributes rather than comparing whole objects.
    \item Search nested collections such as a list of courses, each containing a list of students.
    \item Use a small predicate function or \texttt{lambda} to separate a search procedure from its matching rule.
    \item Explain how \textbf{binary search} uses ordered data to eliminate large portions of the search space.
  \end{itemize}

  \begin{keyidea}
    A search algorithm specifies \emph{where to look next} and \emph{when enough information has been found to stop}.
  \end{keyidea}
\end{frame}

\begin{frame}{The Boundary for This Lecture}
  \begin{columns}[T,onlytextwidth]
    \begin{column}{0.47\textwidth}
      \begin{block}{Today}
        \begin{itemize}
          \item algorithmic problem statements
          \item pseudocode
          \item linear search
          \item searching object attributes
          \item nested searches
          \item simple predicate functions and \texttt{lambda}
          \item binary search
          \item best- and worst-case work
        \end{itemize}
      \end{block}
    \end{column}
    \begin{column}{0.47\textwidth}
      \begin{alertblock}{Not today}
        \begin{itemize}
          \item sorting algorithms
          \item formal correctness proofs
          \item trees or graph search
          \item hash-table implementation
          \item recursive search algorithms
          \item advanced asymptotic analysis
        \end{itemize}
      \end{alertblock}
    \end{column}
  \end{columns}
\end{frame}

\sectionframe{Part 1: From Problems to Algorithms}

\begin{frame}{Searching Is a Computational Problem}
  \begin{columns}[T,onlytextwidth]
    \begin{column}{0.54\textwidth}
      \begin{cpdefinition}
        A \textbf{search problem} asks whether some desired value or condition occurs within a collection and, when useful, where or what the matching item is.
      \end{cpdefinition}

      \begin{itemize}
        \item Is the number \texttt{42} in this list?
        \item At what index does \texttt{"Ada"} occur?
        \item Which student has ID \texttt{1042}?
        \item Does any course contain a student with GPA above \texttt{3.8}?
      \end{itemize}
    \end{column}
    \begin{column}{0.40\textwidth}
      % Intended visual: magnifying glass moving across values, then across Student objects.
      \lectureimage[2.22in]{search-problem-values-and-objects.png}
    \end{column}
  \end{columns}
\end{frame}

\begin{frame}{Definition: Algorithm}
  \begin{cpdefinition}
    An \textbf{algorithm} is a finite, well-defined sequence of steps for solving a problem or performing a computation.
  \end{cpdefinition}

  \begin{itemize}
    \item \textbf{Well-defined:} each step has an unambiguous meaning.
    \item \textbf{Finite:} execution eventually terminates for inputs within the algorithm's intended domain.
    \item \textbf{Effective:} each step can actually be carried out.
    \item \textbf{General:} the procedure describes a family of inputs, not one memorized example.
  \end{itemize}

  \begin{keyidea}
    Python code is one way to express an algorithm; the algorithm itself is the underlying procedure.
  \end{keyidea}
\end{frame}

\begin{frame}{One Problem Can Have More Than One Algorithm}
  Suppose we need to determine whether \texttt{37} appears in a collection.

  \begin{columns}[T,onlytextwidth]
    \begin{column}{0.46\textwidth}
      \begin{block}{Strategy A}
        Start at the beginning and examine each element until the target is found or the collection ends.
      \end{block}
      \begin{center}
        \Large Linear Search
      \end{center}
    \end{column}
    \begin{column}{0.46\textwidth}
      \begin{block}{Strategy B}
        If the values are ordered, examine the middle and repeatedly discard the half that cannot contain the target.
      \end{block}
      \begin{center}
        \Large Binary Search
      \end{center}
    \end{column}
  \end{columns}

  \begin{keyidea}
    The representation and organization of the data determine which algorithms are available.
  \end{keyidea}
\end{frame}

\begin{frame}{Pseudocode Separates the Procedure from Python Syntax}
  \begin{columns}[T,onlytextwidth]
    \begin{column}{0.48\textwidth}
      \begin{block}{Linear-search pseudocode}
        \texttt{LINEAR-SEARCH(values, target)}\\
        \hspace*{0.20in}\texttt{for each position i from left to right}\\
        \hspace*{0.40in}\texttt{if values[i] equals target}\\
        \hspace*{0.60in}\texttt{return i}\\
        \hspace*{0.20in}\texttt{return NOT-FOUND}
      \end{block}
    \end{column}
    \begin{column}{0.46\textwidth}
      \begin{itemize}
        \item Pseudocode communicates the computational structure.
        \item It intentionally omits language-specific punctuation and details.
        \item We can translate the same procedure into Python, Java, C++, or another language.
      \end{itemize}
    \end{column}
  \end{columns}
\end{frame}

\begin{frame}{Correctness, Work, and the Search Contract}
  \begin{columns}[T,onlytextwidth]
    \begin{column}{0.47\textwidth}
      \begin{block}{Before efficiency}
        A useful search must first be:
        \begin{itemize}
          \item \textbf{correct}: returned answers are right;
          \item \textbf{complete}: a match is found when one exists;
          \item \textbf{terminating}: the procedure eventually stops.
        \end{itemize}
      \end{block}
    \end{column}
    \begin{column}{0.47\textwidth}
      \begin{block}{What should the function return?}
        \begin{itemize}
          \item existence $\rightarrow$ \texttt{True}/\texttt{False}
          \item location $\rightarrow$ index
          \item matching item $\rightarrow$ value/object
          \item no match $\rightarrow$ \texttt{-1} or \texttt{None}
        \end{itemize}
      \end{block}
    \end{column}
  \end{columns}

  \begin{keyidea}
    Correctness comes first; then we ask how the work grows as the collection grows.
  \end{keyidea}
\end{frame}

\sectionframe{Part 2: Linear Search}

\begin{frame}{Linear Search Examines Candidates in Sequence}
  \begin{columns}[T,onlytextwidth]
    \begin{column}{0.52\textwidth}
      \begin{cpdefinition}
        \textbf{Linear search} examines elements one at a time until a match is found or every candidate has been checked.
      \end{cpdefinition}

      \begin{itemize}
        \item It does not require the data to be ordered.
        \item It can stop immediately after finding the desired item.
        \item If no match exists, every element must be ruled out.
      \end{itemize}
    \end{column}
    \begin{column}{0.42\textwidth}
      % Intended visual: boxes scanned left-to-right, with each examined box highlighted in sequence.
      \lectureimage[2.22in]{linear-search-left-to-right.png}
    \end{column}
  \end{columns}
\end{frame}

\begin{frame}{Linear Search in Python}
  \begin{block}{Return the index of the first match}
    \texttt{def linear\_search(values: list[int], target: int) -> int:}\\
    \hspace*{0.20in}\texttt{for index, value in enumerate(values):}\\
    \hspace*{0.40in}\texttt{if value == target:}\\
    \hspace*{0.60in}\texttt{return index}\\
    \hspace*{0.20in}\texttt{return -1}
  \end{block}

  \begin{itemize}
    \item Each iteration considers exactly one candidate.
    \item \texttt{return index} terminates the search as soon as the first match is known.
    \item Reaching \texttt{return -1} means the loop exhausted every candidate.
  \end{itemize}
\end{frame}

\begin{frame}{Trace: The Target Position Changes the Work}
  \begin{center}
    \renewcommand{\arraystretch}{1.24}
    \begin{tabular}{l|c|c|l}
      \textbf{call} & \textbf{first match} & \textbf{checks} & \textbf{result} \\\hline
      \texttt{linear\_search([3,7,2,9,5], 7)} & index 1 & 2 & \texttt{1} \\
      \texttt{linear\_search([3,7,2,9,5], 5)} & index 4 & 5 & \texttt{4} \\
      \texttt{linear\_search([3,7,2,9,5], 8)} & absent & 5 & \texttt{-1}
    \end{tabular}
  \end{center}

  \begin{itemize}
    \item A match near the beginning lets the function return early.
    \item A match near the end requires more comparisons.
    \item To conclude ``not found'' in unsorted data, every candidate must be ruled out.
  \end{itemize}
\end{frame}

\begin{frame}{Best Case, Worst Case, and Linear Growth}
  For a list containing $n$ elements:

  \begin{columns}[T,onlytextwidth]
    \begin{column}{0.46\textwidth}
      \begin{block}{Best case}
        The first element matches: approximately one comparison.
      \end{block}
    \end{column}
    \begin{column}{0.46\textwidth}
      \begin{block}{Worst case}
        The target is last or absent: approximately $n$ comparisons.
      \end{block}
    \end{column}
  \end{columns}

  \begin{keyidea}
    Worst-case linear search grows as $O(n)$: doubling the number of candidates can approximately double the required checks.
  \end{keyidea}
\end{frame}

\begin{frame}{Linear Search Is More General Than Equality}
  The search procedure can test any condition that identifies the desired item.

  \begin{block}{Examples of matching rules}
    \begin{itemize}
      \item \texttt{value == 42}
      \item \texttt{name == "Ada"}
      \item \texttt{student.student\_id == 1042}
      \item \texttt{student.gpa >= 3.8}
      \item \texttt{course.code == "CS1001"}
    \end{itemize}
  \end{block}

  \begin{keyidea}
    Searching means examining candidates until one satisfies the condition that defines a match.
  \end{keyidea}
\end{frame}

\sectionframe{Part 3: Searching Collections of Objects}

\begin{frame}{Objects Change What We Compare}
  \begin{block}{Student class}
    \texttt{class Student:}\\
    \hspace*{0.20in}\texttt{def \_\_init\_\_(self, student\_id: int, name: str, gpa: float):}\\
    \hspace*{0.40in}\texttt{self.student\_id = student\_id}\\
    \hspace*{0.40in}\texttt{self.name = name}\\
    \hspace*{0.40in}\texttt{self.gpa = gpa}
  \end{block}

  \begin{itemize}
    \item A list can store references to \texttt{Student} instances.
    \item The target may not be an entire object.
    \item Usually we search by one meaningful attribute or condition.
  \end{itemize}
\end{frame}

\begin{frame}{Search a List of Objects by an Attribute}
  \begin{block}{Return the matching student}
    \texttt{def find\_student(students: list[Student], target\_id: int):}\\
    \hspace*{0.20in}\texttt{for student in students:}\\
    \hspace*{0.40in}\texttt{if student.student\_id == target\_id:}\\
    \hspace*{0.60in}\texttt{return student}\\
    \hspace*{0.20in}\texttt{return None}
  \end{block}

  \begin{keyidea}
    The algorithm is still linear search. Only the matching rule changed from comparing primitive values to inspecting an object's attribute.
  \end{keyidea}
\end{frame}

\begin{frame}{Trace: Object Search}
  \begin{columns}[T,onlytextwidth]
    \begin{column}{0.55\textwidth}
      \begin{block}{Students}
        \texttt{[Student(1008, "Maya", 3.4),}\\
        \texttt{ Student(1042, "Luis", 3.8),}\\
        \texttt{ Student(1091, "Nora", 3.6)]}
      \end{block}

      Search for ID \texttt{1042}:
      \begin{enumerate}
        \item inspect the object at position 0;
        \item read \texttt{student.student\_id};
        \item no match, advance;
        \item inspect position 1;
        \item matching attribute found, return that object.
      \end{enumerate}
    \end{column}
    \begin{column}{0.38\textwidth}
      % Intended visual: list slots pointing to Student objects, with student_id field highlighted during search.
      \lectureimage[2.22in]{linear-search-student-objects.png}
    \end{column}
  \end{columns}
\end{frame}

\begin{frame}{What Should an Object Search Return?}
  Suppose a matching \texttt{Student} is found.

  \begin{center}
    \renewcommand{\arraystretch}{1.30}
    \begin{tabular}{l|l}
      \textbf{Return} & \textbf{Useful when the caller needs...} \\\hline
      \texttt{True} & only whether a match exists \\
      index & the object's position in the list \\
      \texttt{Student} object & access to the matching student's state or methods \\
      \texttt{None} & a clear ``no matching object'' result
    \end{tabular}
  \end{center}

  \begin{keyidea}
    Returning the matching object is often more useful than returning its index because the caller can immediately access its attributes and methods.
  \end{keyidea}
\end{frame}

\begin{frame}{From Specific Searches to a General Predicate}
  Exact IDs are only one matching rule. A search can instead receive a function that answers whether one candidate matches.

  \begin{block}{Reusable search}
    \texttt{def find\_first(items, matches):}\\
    \hspace*{0.20in}\texttt{for item in items:}\\
    \hspace*{0.40in}\texttt{if matches(item):}\\
    \hspace*{0.60in}\texttt{return item}\\
    \hspace*{0.20in}\texttt{return None}
  \end{block}

  \begin{itemize}
    \item \texttt{find\_first} controls \textbf{how} candidates are examined.
    \item \texttt{matches} determines \textbf{which} candidates qualify.
  \end{itemize}
\end{frame}

\begin{frame}{A Small Use for \texttt{lambda}}
  A \texttt{lambda} expression creates a compact function value for a one-use matching rule.

  \begin{block}{Same search procedure, different predicates}
    \texttt{find\_first(students, lambda s: s.student\_id == 1042)}\\[0.05in]
    \texttt{find\_first(students, lambda s: s.name == "Maya")}\\[0.05in]
    \texttt{find\_first(students, lambda s: s.gpa >= 3.8)}
  \end{block}

  \begin{keyidea}
    The \texttt{lambda} is not a new search algorithm. It only supplies the condition used by the existing linear search.
  \end{keyidea}
\end{frame}

\sectionframe{Part 4: Searching Nested Collections}

\begin{frame}{Objects Can Contain Collections of Other Objects}
  \begin{block}{Course class}
    \texttt{class Course:}\\
    \hspace*{0.20in}\texttt{def \_\_init\_\_(self, code: str, students: list[Student]):}\\
    \hspace*{0.40in}\texttt{self.code = code}\\
    \hspace*{0.40in}\texttt{self.students = students}
  \end{block}

  A program might then store:
  \begin{block}{Outer collection}
    \texttt{courses = [course1, course2, course3]}
  \end{block}

  \begin{center}
    \texttt{list[Course]} $\rightarrow$ each \texttt{Course} $\rightarrow$ \texttt{list[Student]}
  \end{center}
\end{frame}

\begin{frame}{The Search Must Follow the Shape of the Data}
  \begin{columns}[T,onlytextwidth]
    \begin{column}{0.53\textwidth}
      To search every student across every course:
      \begin{enumerate}
        \item examine each \texttt{Course} in the outer list;
        \item enter that course's \texttt{students} list;
        \item examine each \texttt{Student};
        \item stop as soon as the matching student is found.
      \end{enumerate}

      \begin{keyidea}
        Nested structure naturally produces nested traversal.
      \end{keyidea}
    \end{column}
    \begin{column}{0.40\textwidth}
      % Intended visual: outer list of Course objects, each Course pointing to an inner list of Student objects.
      \lectureimage[2.05in]{nested-course-student-search.png}
    \end{column}
  \end{columns}
  {\footnotesize Student instances are labeled by ID; the list slots store references.}
\end{frame}

\begin{frame}{Search a List of Objects with Nested Lists of Objects}
  \begin{block}{Return the first matching student}
    \texttt{def find\_student\_in\_courses(courses: list[Course], target\_id: int):}\\
    \hspace*{0.20in}\texttt{for course in courses:}\\
    \hspace*{0.40in}\texttt{for student in course.students:}\\
    \hspace*{0.60in}\texttt{if student.student\_id == target\_id:}\\
    \hspace*{0.80in}\texttt{return student}\\
    \hspace*{0.20in}\texttt{return None}
  \end{block}

  \begin{itemize}
    \item The outer loop moves among course objects.
    \item The inner loop moves among student objects belonging to the current course.
    \item A successful inner comparison terminates the entire function.
  \end{itemize}
\end{frame}

\begin{frame}{Trace the Nested Search --- and Preserve Context}
  Suppose \texttt{CS1001} contains IDs \texttt{[1008, 1014]} and \texttt{CS2020} contains \texttt{[1042, 1077]}.

  \begin{center}
    \renewcommand{\arraystretch}{1.20}
    \begin{tabular}{l|l|l}
      \textbf{course} & \textbf{student ID} & \textbf{result} \\\hline
      CS1001 & 1008 & continue \\
      CS1001 & 1014 & continue \\
      CS2020 & 1042 & match
    \end{tabular}
  \end{center}

  Returning \texttt{student} gives the matching object. Returning \texttt{(course, student)} preserves both the match and the container in which it was found.

  \begin{keyidea}
    Search design includes deciding what context should accompany a successful match.
  \end{keyidea}
\end{frame}

\begin{frame}{How Much Work Does the Nested Search Perform?}
  Suppose the courses contain \texttt{4}, \texttt{6}, and \texttt{5} students.

  \begin{itemize}
    \item A worst-case search can inspect all \texttt{4 + 6 + 5 = 15} students.
    \item Count \textbf{all student entries} plus visits to the course containers.
    \item The nested loops do not automatically mean the algorithm is ``quadratic.''
  \end{itemize}

  \begin{block}{General rule}
    Total worst-case work is $O(C+S)$: $C$ is the number of courses and $S$ is the total number of student entries. Even an empty course requires an outer-loop visit.
  \end{block}
\end{frame}

\begin{frame}{A Common Nested-Search Bug: Returning Too Early}
  \begin{columns}[T,onlytextwidth]
    \begin{column}{0.47\textwidth}
      \begin{alertblock}{Incorrect}
        \footnotesize
        \texttt{for course in courses:}\\
        \hspace*{0.20in}\texttt{for student in course.students:}\\
        \hspace*{0.40in}\texttt{if student.student\_id == target\_id:}\\
        \hspace*{0.60in}\texttt{return student}\\
        \hspace*{0.40in}\texttt{return None}
      \end{alertblock}
    \end{column}
    \begin{column}{0.47\textwidth}
      \begin{block}{Problem}
        The first nonmatching student immediately returns \texttt{None}. The function never examines later candidates.
      \end{block}

      \begin{keyidea}
        ``Not found'' is justified only after every relevant candidate has been ruled out.
      \end{keyidea}
    \end{column}
  \end{columns}
\end{frame}

\sectionframe{Part 5: Binary Search}

\begin{frame}{Linear Search Does Not Use Ordering Information}
  Consider an ordered list:

  \begin{center}
    \Large\texttt{[3, 7, 12, 18, 24, 31, 42, 50, 61]}
  \end{center}

  Searching for \texttt{42} linearly still begins at \texttt{3}, then \texttt{7}, then \texttt{12}, and so on.

  \begin{block}{Observation}
    Because the values are ordered, comparing the target with one carefully chosen element can rule out many positions at once.
  \end{block}
\end{frame}

\begin{frame}{Binary Search Requires Ordered Data}
  \begin{cpdefinition}
    \textbf{Binary search} repeatedly examines the middle of an ordered search interval and discards the half that cannot contain the target.
  \end{cpdefinition}

  \begin{itemize}
    \item If the middle value equals the target, stop.
    \item If the target is smaller, continue in the left half.
    \item If the target is larger, continue in the right half.
  \end{itemize}

  \begin{alertblock}{Precondition}
    Binary search depends on the collection being ordered in ascending order by the property being compared.
  \end{alertblock}
\end{frame}

\begin{frame}{Binary Search Shrinks the Search Interval}
  Search for \texttt{42}: compare with \texttt{24}, keep indices 5--8, then inspect index 6.
  \begin{center}
    \lectureimage[1.90in]{binary-search-halving-interval.png}
  \end{center}
\end{frame}

\begin{frame}{Binary Search Maintains a Shrinking Interval}
  \small
  Precondition: \texttt{values} is sorted in ascending order.
  \begin{block}{Implementation}
    \texttt{def binary\_search(values: list[int], target: int) -> int:}\\
    \hspace*{0.20in}\texttt{left = 0}\\
    \hspace*{0.20in}\texttt{right = len(values) - 1}\\
    \hspace*{0.20in}\texttt{while left <= right:}\\
    \hspace*{0.40in}\texttt{mid = (left + right) // 2}\\
    \hspace*{0.40in}\texttt{if values[mid] == target: return mid}\\
    \hspace*{0.40in}\texttt{elif values[mid] < target: left = mid + 1}\\
    \hspace*{0.40in}\texttt{else: right = mid - 1}\\
    \hspace*{0.20in}\texttt{return -1}
  \end{block}

  Each comparison finds a match or shrinks \texttt{left ... right}. With duplicate values, this version may return any matching index.
\end{frame}

\begin{frame}{Trace: Binary Search for \texttt{31}}
  Search \texttt{[3, 7, 12, 18, 24, 31, 42, 50, 61]}.

  \begin{center}
    \renewcommand{\arraystretch}{1.23}
    \begin{tabular}{c|c|c|c|l}
      \textbf{left} & \textbf{right} & \textbf{mid} & \textbf{value} & \textbf{action} \\\hline
      0 & 8 & 4 & 24 & move \texttt{left} to 5 \\
      5 & 8 & 6 & 42 & move \texttt{right} to 5 \\
      5 & 5 & 5 & 31 & return 5
    \end{tabular}
  \end{center}

  Only three elements were inspected.
\end{frame}

\begin{frame}{Why Binary Search Is Faster on Large Ordered Data}
  \small
  Each unsuccessful comparison removes approximately half of the remaining candidates.

  \begin{center}
    \renewcommand{\arraystretch}{1.20}
    \begin{tabular}{c|c}
      \textbf{number of candidates} & \textbf{about how many halvings?} \\\hline
      8 & 3 \\
      16 & 4 \\
      1,024 & 10 \\
      1,048,576 & 20
    \end{tabular}
  \end{center}

  \begin{keyidea}
    For $n\ge1$ candidates, this implementation inspects at most $\lfloor\log_2 n\rfloor+1$ middle elements. Halvings and inspected elements differ by up to one; growth is $O(\log n)$.
  \end{keyidea}
\end{frame}

\begin{frame}{Binary Search Fails Without the Ordering Precondition}
  Consider:
  \begin{center}
    \Large\texttt{[18, 3, 42, 7, 31, 12, 24]}
  \end{center}

  If the middle value is \texttt{7} and the target is \texttt{31}, we cannot conclude that \texttt{31} must lie to the right.

  \begin{alertblock}{Why not?}
    Values smaller and larger than \texttt{7} appear on both sides. The comparison does not identify a half that can safely be discarded.
  \end{alertblock}

  \begin{keyidea}
    Binary search is powerful because ordering turns one comparison into information about many positions.
  \end{keyidea}
\end{frame}

\begin{frame}{Binary Search Can Search Objects --- by the Ordering Attribute}
  Suppose the list is ordered by \texttt{student\_id}.

  \begin{block}{Middle comparison}
    \texttt{student = students[mid]}\\
    \texttt{if student.student\_id == target\_id: ...}\\
    \texttt{elif student.student\_id < target\_id: ...}
  \end{block}

  \begin{itemize}
    \item We compare the attribute that defines the ordering, not the entire object.
    \item A list ordered by ID supports binary search by ID.
    \item That same list is \emph{not} automatically suitable for binary search by GPA or name.
  \end{itemize}

  \begin{keyidea}
    The ordering and the search comparison must refer to the same property.
  \end{keyidea}
\end{frame}

\begin{frame}{Linear Search vs. Binary Search}
  \begin{center}
    \renewcommand{\arraystretch}{1.32}
    \begin{tabular}{l|c|c}
      & \textbf{Linear Search} & \textbf{Binary Search} \\\hline
      Requires ordered data? & No & Yes \\
      Candidate examined next & next element & middle of interval \\
      Worst-case growth & $O(n)$ & $O(\log n)$ \\
      Works naturally on arbitrary predicates? & Yes & only with compatible ordering \\
      Simple nested object search? & Yes & usually not directly
    \end{tabular}
  \end{center}
\end{frame}

\begin{frame}{Representation Determines the Available Search}
  \begin{center}
    \renewcommand{\arraystretch}{1.25}
    \begin{tabular}{l|l}
      \textbf{representation} & \textbf{natural search approach} \\\hline
      unordered list & linear search \\
      ordered list & linear or binary search \\
      dictionary by key & dictionary key lookup \\
      nested object collections & traversal following the nested structure
    \end{tabular}
  \end{center}

  \begin{block}{Make assumptions visible}
    A binary-search function should clearly communicate its precondition: the collection must already be ordered by the property being compared.
  \end{block}

  \begin{keyidea}
    Choosing a data representation and choosing an algorithm are connected design decisions.
  \end{keyidea}
\end{frame}

\begin{frame}{Integrated Trace: Which Search Is Appropriate?}
  For each task, choose the most natural search strategy.

  \begin{enumerate}
    \item Find the first student with GPA above \texttt{3.8} in an unsorted list.
    \item Find student ID \texttt{1042} in a list already ordered by student ID.
    \item Find student ID \texttt{1042} somewhere inside a list of course objects, each containing its own student list.
    \item Retrieve the value associated with key \texttt{"Ada"} from a dictionary.
  \end{enumerate}

  \begin{block}{Reason from structure}
    Ask what organization the data already provides and what condition defines a match.
  \end{block}
\end{frame}

\begin{frame}{Takeaways}
  \begin{enumerate}
    \item An algorithm is a precise, finite procedure for solving a computational problem.
    \item Linear search checks candidates sequentially and works even when the data is unordered.
    \item Searching objects usually means examining a selected attribute or predicate.
    \item Nested object structures require traversal that follows the nesting of the data.
    \item A \texttt{lambda} can compactly provide a one-use matching rule to a reusable search function.
    \item Binary search repeatedly discards half of an \textbf{ordered} search interval.
    \item Linear search has worst-case growth $O(n)$; binary search has worst-case growth $O(\log n)$.
    \item Data organization determines which search strategies are valid and useful.
  \end{enumerate}
\end{frame}

\end{document}
